\section{Introduction}\label{sec:intro}

Language models can now take part in research well beyond answering questions.  They write and run numerical code, produce formal
proofs, search the literature and draft text, over sessions that add up to months.  For a single researcher without institutional
support this changes what is feasible.  It also changes what can go wrong.  A long record written in dialogue with a model accumulates
errors of a particular kind, and they are hard to see because each step reads plausibly in prose.  In the project described in
Section~\ref{sec:case} three such classes dominated:
\begin{enumerate}
\item \emph{Frame mixing.}  A quantity defined in one reference frame, or of one kind (a scale, a readout), is compared with a quantity
from another as if they were the same.
\item \emph{Input read as output.}  A value that was put in, or chosen to fit, is later cited as if it had been derived.
\item \emph{Promotion.}  A local, conditional or numerical statement is restated as a theorem, and later steps build on the stronger form.
\end{enumerate}
Peer review would catch some of these, but a long record produced at the speed of a model is not reviewed step by step, and the model
itself is not a reliable reviewer of its own earlier output.

This paper describes a pragmatic answer: a small typed language, QAM, in which the record is kept, and a store machine, MMA, which
executes instructions on the record and logs every step.  The language does not try to formalize the mathematics or the physics.  It
types the \emph{objects of the record} --- states, layers, operators, readouts, branches, shells, witnesses and tensor composites ---
and makes the three error classes ill-formed: a comparison across frames needs a declared comparison object, a value that rests on an
input carries a witness, and a certificate for a statement can only be formed from a proof.  The core of the language and of the machine
is formalized in Lean~4~\cite{lean4}.

\paragraph{Contributions.}
\begin{enumerate}
\item The language QAM over Zermelo--Fraenkel set theory, with its two-sided display for frame-typed quantities, and the instruction
layer MML with its transition semantics MMA (Sections~\ref{sec:language} and~\ref{sec:machine}).
\item Lean proofs of the properties that make the language usable as a record: disjointness of object kinds, an exactly invertible
reindexing, preservation of closure classes along admissible transport, and for the machine preservation, progress for the
non-modifying instructions, atomicity, ledger completeness and the absence of false certificates; and a project-independent module that
states the three error classes as rules on derivation trees and proves that well-formed trees never promote, never lose or invent an
input, and never change frame except through declared maps (Section~\ref{sec:lean}).
\item A \LaTeX{} package, \texttt{qam.sty}, that applies the same rules while a document is compiled and stops the compiler at the first
violation, that generates the QAM notation of each claim from its declaration, that links every claim to its sources by a chain of hashes starting at the thesis, and that exports
the record to a file which a verified Lean checker reads (Section~\ref{sec:tex}).
\item A case study of a research project in mathematical physics, carried out by one author with language models from the beginning,
with about a hundred append-only audit rounds and more than sixty pre-registered tests (Section~\ref{sec:case}).
\item An account of what the discipline caught and what it did not, and of what the formal guarantees mean (Section~\ref{sec:limits}).
\end{enumerate}

The guarantees are about form.  QAM can ensure that a record never compares across frames without saying so, never cites an input as a
derivation, and never certifies what was not proved.  It cannot ensure that a premise is true, that a literature value was read
correctly, or that a physical reading makes sense.  Section~\ref{sec:case} contains an example of each.
